\documentclass[11pt]{article}
\usepackage[a4paper,margin=2.5cm]{geometry}
\usepackage{amsmath,amssymb}
\usepackage{graphicx}
\usepackage{booktabs}
\usepackage{parskip}

\title{MA-INF 2201 Computer Vision WS26/27 \\ Exercise Sheet 3 --- Written Answers}
\author{Mohamadhosein Sabaghi \and Sepehr Noey Kordkandi}
\date{}

\begin{document}
\maketitle

% ==========================================================================
\section{Hough Transform for Line Detection}

Edge points used: the 5 bright pixels of \texttt{control.png} at
$(8,64), (64,8), (64,64), (64,120), (120,64)$ --- a plus-sign arrangement with an
obvious horizontal line through $(8,64),(64,64),(120,64)$ and an obvious vertical
line through $(64,8),(64,64),(64,120)$.

\subsection{Why the Hough Space Looks Like Sine Waves}

For a fixed point $(x,y)$, the normal-form line equation gives
\[
\rho(\theta) = x\cos\theta + y\sin\theta.
\]
This is exactly the trigonometric identity $a\cos\theta + b\sin\theta = R\cos(\theta-\varphi)$
with $R=\sqrt{x^2+y^2}$ and $\varphi=\operatorname{atan2}(y,x)$ --- i.e.\ a sinusoid in
$\theta$ with amplitude equal to the point's distance from the origin and a phase shift
equal to the point's angular position. Every point in the image therefore traces out one
sine curve in $(\theta,\rho)$ space, and points that are collinear in the image produce
curves that all cross at the same $(\theta,\rho)$ --- that crossing is exactly the line's
own normal-form parameters, which is why the accumulator bin there collects one vote per
point on the line.

\subsection{Can the Accumulator Have a Straight Horizontal Line?}

Yes. A horizontal line in Hough space means $\rho$ stays constant while $\theta$ varies
over the full range, i.e.\ $x\cos\theta + y\sin\theta$ is the same value for every $\theta$.
Since $\cos\theta$ and $\sin\theta$ vary independently as $\theta$ sweeps its range, the only
way their weighted sum can be constant for all $\theta$ is if $x=0$ and $y=0$: that is, a point
located exactly at the origin $(0,0)$ produces a flat line $\rho(\theta)=0$ for every $\theta$.
None of the 5 points in \texttt{control.png} sit at the origin (the origin is the image's
top-left corner), so no such horizontal line appears in our accumulator --- all 5 points lie
strictly inside the image, away from $(0,0)$.

\subsection{$\theta\in[0^\circ,360^\circ)$ and Multiple Representations}

With the default range $\theta\in[50^\circ,270^\circ]$, thresholding the accumulator at
$\geq 3$ votes already gives \textbf{three} detected lines, not two:

\begin{center}
\begin{tabular}{ccc}
\toprule
$\theta$ & $\rho$ & votes \\
\midrule
$90^\circ$ & $64$ & 3 \\
$270^\circ$ & $-64$ & 3 \\
$180^\circ$ & $-64$ & 3 \\
\bottomrule
\end{tabular}
\end{center}

Checking what these actually correspond to: $(\theta,\rho)=(90^\circ,64)$ gives
$0\cdot x + 1\cdot y = 64$, i.e.\ the line $y=64$; $(\theta,\rho)=(270^\circ,-64)$ gives
$0\cdot x + (-1)\cdot y = -64$, which is \emph{also} $y=64$ --- the same horizontal line, in
its second valid parametrization. $(\theta,\rho)=(180^\circ,-64)$ gives $-1\cdot x = -64$,
i.e.\ $x=64$, the vertical line.

Every line has exactly two equivalent normal-form representations, $(\theta,\rho)$ and
$(\theta+180^\circ,-\rho)$, since reversing the normal direction and negating the signed
distance describes the same geometric line. The default range already spans $220^\circ$
(more than the $180^\circ$ needed for a unique representation), so it happens to capture
\emph{both} representations of the horizontal line ($90^\circ$ and $270^\circ$ are
$180^\circ$ apart and both lie in $[50^\circ,270^\circ]$), while only capturing one
representation of the vertical line (its other form, $(0^\circ,64)$, falls just outside the
range since $\mathrm{THETA\_MIN}=50^\circ$).

Extending the range to the full $[0^\circ,360^\circ)$ brings every line's \emph{second}
representation into range as well --- here, specifically $(\theta,\rho)=(0^\circ,64)$ for
the vertical line. So instead of finding 2 physical lines, we expect to detect
\textbf{4} accumulator peaks for the same 2 lines: this is why there end up being more than
two representations overall once the range exceeds $180^\circ$ --- each physical line
contributes one extra, numerically distinct but geometrically identical, peak for every
full $180^\circ$ of extra range covered.
% TODO: confirm this matches the actual printed output after rerunning q1_template.py
% with THETA_MIN=0, THETA_MAX=360, and adjust the exact count/values if needed.

% ==========================================================================
\section{Roll Correction with Hough Transform}
% TODO: fill in once q2_template.py has been run.
% - Visualise/describe the Canny edges, top-10 detected lines, and the unrolled image.
% - Explain how adding more lines affects the correction angle (e.g. averaging
%   more line-angle estimates reduces sensitivity to any single noisy/outlier
%   detected line, but may also pull in lines that are not aligned with the
%   dominant roll direction if the threshold is too permissive).

TODO: write this section.

% ==========================================================================
\section{Hough Transform for Circle Detection}
% TODO: fill in once q3_template.py has been run.
% - Describe the detected circles overlaid on coins.jpg.
% - Describe the accumulator slices at different radii and at the best radius.
% - Comment on how results vary with threshold and radius range (e.g. too low
%   a threshold gives false positives, too narrow a radius range misses coins
%   of a different size, too high a threshold misses faint/partial circles).

TODO: write this section.

\end{document}
